\section{Motivation}

The rapid growth of machine learning model complexity from large language models to diffusion-based generative architectures has created a widening gap between the compute resources available to industry and those accessible to academic researchers. While large technology companies operate dedicated GPU clusters with thousands of accelerators, university departments typically manage a handful of machines with far more modest capabilities. This disparity forces academic AI research groups to either compete for shared institutional HPC resources with long queuing times, or to maintain their own small-scale infrastructure with limited administrative support.

A common response to this challenge in mid-size research institutions is to adopt two separate infrastructure stacks: a traditional HPC scheduler such as Slurm for batch training jobs, and a cloud-native platform such as Kubernetes for long-running interactive services like JupyterHub and experiment tracking dashboards. This \emph{infrastructure sprawl} introduces significant operational overhead, as both stacks require dedicated DevOps expertise to deploy, configure, and maintain. For small academic departments that cannot afford a dedicated infrastructure team, this dual-stack approach is often impractical, leaving researchers to fall back on ad-hoc, unmanaged access to individual machines.

\subsection{Computing at our department}

On a global scale, Comenius University is a rather small institution, with our faculty being an even smaller one, let alone our department. The Department of Applied Informatics has around \todo{Insert exact number of faculty members.} faculty members and around \todo{Insert exact number of students.} students. Many of these students are actively doing research in machine learning and adjacent fields, thus requiring computational resources that often go beyond their own personal computers.

The department operates a heterogeneous collection of compute nodes built from consumer-grade hardware. The machines are equipped with a mix of GPUs spanning several generations---including NVIDIA Titan V, Titan XP, GeForce RTX 3060, and GeForce RTX 4090 with VRAM capacities ranging from 12\,GB to 24\,GB. CPU configurations and RAM sizes similarly vary across nodes. This hardware heterogeneity, while cost-effective, presents its own set of challenges for workload placement and resource management. On request, these machines are available to any student who needs them for their work.

\subsection{Current state}

Currently there are four machines used for compute. Users are provided access to each machine the department manages based on a request from their supervisors. They are also added to a Teams team, which is used to post updates about the machines. It is also used for ``reserving'' compute slots users post to the team when, for how long, and on which machine they are going to be computing. Users then use SSH to connect to the selected nodes, move their data to the node, and run their computations.

However, because the computers are standalone machines with no collective management or shared file system beyond user logins, this setup comes with a number of challenges: environment management, data management, and scheduling. These are discussed in detail in the problem statement.

This ad-hoc arrangement illustrates the broader problem facing small research groups: without a unified platform, departments must choose between operating without any orchestration accepting the inefficiencies described above or adopting heavyweight infrastructure stacks that require dedicated administrators. Neither extreme serves the needs of a small academic lab well. The solution explored in this thesis seeks a middle ground: a \emph{single}, lightweight control plane that consolidates batch workloads, interactive development environments, and supporting services onto one platform manageable by a single administrator.
